\section{Removing the disk obstruction}\label{sec:curvature}

Let $i:L\to D$ be the graded spin immersion of
Proposition~\ref{prop:immersion}.  Its transverse self-intersection
generators have odd integer degrees, and
\begin{equation}\label{curv:topological-input}
 H^2(L,\Q)\simeq H^2(D,\Q)/\mathcal Z,
 \qquad
 i_*:H_{n-2}(L,\Q)\longrightarrow H_{n-2}(D,\Q)
 \text{ is injective}.
\end{equation}
We will use the second assertion to remove the curvature of its
Fukaya algebra.  The first will then bound its degree-two endomorphisms.
The geometric ingredient is a Stokes identity for disks with one
interior marked point and no external boundary marked point.  We
construct that identity geometrically, then use filtered transfer to
remove curvature.

\subsection{Branes, gradings, and the filtered construction}

Choose countably many Hamiltonian copies of the product test branes,
ordered by the positive integers, with each test occurring cofinally.
Use product Hamiltonian perturbations, small enough that the torus
factors remain perturbed graphs.  Arrange pairwise transversality and
absence of triple intersections, including intersections with $L$.
When both factors are present, impose these conditions separately on
the factors of all embedded copies.  Hamiltonian variations supply the
required tangent directions, so these are countably many generic
conditions.

Fix a split compatible almost complex structure $J_0$ with no
nonconstant spheres and no nonconstant disks on any embedded test
copy.  The sphere and disk exclusion on the quartic factor is the
genericity assertion used in the test-brane construction; on the torus
factor it follows from the vanishing of the relevant second relative
homotopy groups.  No disk exclusion is imposed on $L$.  All category
constructions in this section involve a finite subcollection of the
copies together with $L$.

Write $\Lambda_0$ for the universal Novikov ring with complex
coefficients and $\Lambda$ for its fraction field; its monomials are
written $T^\lambda$, with valuation $\lambda$.  We use the immersed
cyclic de Rham construction of Abouzaid--Fukaya--Oh--Ohta--Ono
\cite[Theorems~9.3 and~9.6]{AFOOO2026}.  It gives a curved cyclic
filtered $A_\infty$ category with homotopy units for finite collections
of closed spin immersions in a compact symplectic manifold.  Here the
relative spin background and bulk class are zero.  The morphism spaces
are forms on the source diagonals and the generators at transverse
intersections.  We denote the operations by $m_k$, their cyclic pairing
by $\langle\ ,\ \rangle$, and the curvature by $m_0$.

These operations have integer degree $2-k$.  Indeed, the brane
gradings specify integer indices for the cap paths at every transverse
end, including self-intersection ends.  After capping, the boundary
grading has Maslov index zero.  There is no ambiguity from spheres
because the ambient first Chern class vanishes.  The same index
calculation applies to parameterized operations.  Energy monomials
have degree zero, and all contractions, transfers, and comparison
maps below are chosen homogeneous.  Thus the grading is retained by
the construction itself, including the shifted degrees or auxiliary
variables used in a particular model.

We use Kuranishi structures with normalized boundaries and corners,
outer collars where required, and continuous-family perturbations
(CF perturbations).  The local chart theorem and its evaluation
submersion criterion are
\cite[Theorem~7.1, Condition~7.11, and Lemma~7.12]{FOOOChartsI};
compatible systems arise from disk-component-wise obstruction data
\cite[Theorems~5.3 and~5.4]{FOOOChartsII}.  We use the corrected chart
conventions of \cite{FOOOChartsICorrigendum}.
The analytic gluing input is the exponential control of the gluing
maps and their length derivatives
\cite[Assumption~3.12 and Theorems~3.13, 6.4, 8.16, and~8.17]{FOOOGluing}.
For transverse immersed ends, the fixed-jump charts and spin
determinant cappings are those of
\cite[Sections~4--5]{AkahoJoyce}, as incorporated in the immersed
construction of \cite{AFOOO2026}.

In particular, the nodal chart uses component surjectivity augmented
by matching equations at ordinary boundary nodes and by the transverse
asymptotic conditions at jumps.  Auxiliary markings stabilize the
domains, and the obstruction spaces are transported on compact cores
before gluing.  These are the local inputs needed below when the total
curve has no external boundary mark.  All the additional spaces used
for that purpose have one interior mark and no sphere
components; there is no gluing at an angular interior node.

The construction is first made through finite energy and arity
cutoffs.  Filtered pseudo-isotopies, homotopy limits, and canonical
models then give the complete category
\cite[Sections~6, 9--10, and~12, in particular Propositions~9.55
and~12.10]{AFOOO2026}.  We keep $J_0$ fixed in this passage.  The
following observation gives a common energy filtration for a finite
list of branes.

\begin{lemma}[Energy values]\label{curv:energy}
For a finite list of the immersed and embedded branes above, the
energies of $J_0$-holomorphic polygons form a locally finite subset of
$\R_{\geq0}$.  In particular, the nonconstant polygons have a common
positive energy lower bound, independent of their number of marked
points.
\end{lemma}

\begin{proof}
Let $K\subset D$ be the union of their images.  It has the homotopy
type of a finite complex and has arbitrarily small regular
neighborhoods: away from the finitely many double points these are
ordinary tubular neighborhoods, and at a double point they are the
standard transverse-double neighborhoods.  Relative de Rham theory
therefore represents every class in $H^2(D,K;\R)$ by a smooth closed
form that vanishes near $K$.

The symplectic form has a relative lift that computes the polygon
areas.  To see this directly, construct near $K$ a primitive $\alpha$
of $\omega$ whose pullback to every sheet is zero.  At a double point,
simultaneous straightening of the transverse Lagrangians gives a radial
primitive with this property.  On the remaining pieces use the
Lagrangian neighborhood theorem.  After shrinking the neighborhoods,
their overlaps at the ends of the pieces are collars on disjoint
sheets.  The differences of the primitives are exact there, with
potentials zero on the sheets, by the relative Poincar\'e lemma.
These potentials patch the primitives.  Cutting off $\alpha$ gives a
closed form $\widetilde\omega$ vanishing near $K$ and representing the
relative lift.  Stokes' theorem, with the usual decay at jump ends,
gives
\[
 \int u^*\widetilde\omega=\int u^*\omega
\]
for every polygon $u$.

Choose closed forms vanishing near $K$ that represent a basis of
$H^2(D,K;\R)$.  Their integrals over a $J_0$-holomorphic polygon are
bounded in absolute value by fixed constants times its area: this is
the pointwise comass estimate for the compatible metric.  A polygon
defines an integral class in $H_2(D,K;\Z)$, whose quotient by torsion
is a lattice.  The preceding bounds leave only finitely many lattice
points below a fixed area cutoff.  Pairing with
$[\widetilde\omega]$ leaves only finitely many energies.  Energy
positivity now gives the asserted lower bound whenever a nonconstant
polygon exists.  If none exists, any positive lower bound will do.
\end{proof}

The additive monoid generated by these energies is locally finite as
well: below an energy bound $B$, each summand belongs to a finite set
and a sum has at most $B/\delta$ positive summands, where $\delta>0$
is the gap.  This is the filtration used for the transfer operations.

At energy zero, the unary operation is the de Rham differential on
the source diagonals and zero on transverse-point generators; the
constant strip with no stabilizing mark is unstable.  Transfer at
energy zero consequently gives the graded canonical space
\begin{equation}\label{curv:canonical-space}
 C_L^{\mathrm{can}}
 =H^*(L;\C)\otimes\Lambda_0
   \;\oplus\!\!\bigoplus_{x\text{ a branch generator}}\!\Lambda_0 x.
\end{equation}
The diagonal transfer can be made cyclic using a compatible de Rham
contraction.  Homotopy units will give cohomological units after
curvature is removed.  Since every branch generator has odd degree,
the degree-two part of this canonical space is exactly
$H^2(L;\C)\otimes\Lambda_0$.  Removing curvature will therefore
give the cohomological bound in the following proposition, which is
the output of this section.

\begin{proposition}\label{prop:curvature}
For every finite subcollection of the embedded test copies, $L$ and
those copies define an uncurved proper cohomologically unital cyclic
$A_\infty$ category over $\Lambda$.  Its split-closed triangulated
hull has Serre functor $[n]$, and
\begin{equation}\label{curv:degree-two-bound}
 \dim_\Lambda \Hom^2(L,L)
 \leq b_2(L)=b_2(D)-\dim_\Q\mathcal Z.
\end{equation}
For a sequence of embedded copies with strictly increasing labels,
the operations are the geometric disk operations at $J_0$, up to
filtered pseudo-isotopy through every finite cutoff.
\end{proposition}

The proof occupies the rest of the section.  To remove curvature, we construct scalar operations on disks with
an interior root, relative to the ordinary cyclic disk operations.
Their Stokes identity detects the leading curvature class against
ambient cohomology and remains valid after positive-valuation
boundary insertions.  Filtered transfer then gives the proposition.

\subsection{A Stokes identity with one interior mark}

Fix the collared componentwise Kuranishi data and CF-perturbations
defining the cyclic boundary-rooted disk operations through a finite
energy and arity cutoff, enlarged to contain the factors of all
boundary strata that will occur.  We construct additional scalar
operations relative to these geometric choices.  Their disks carry one distinguished
interior mark and cyclically ordered boundary inputs, but no boundary
output.  The component containing the interior mark roots a broken
disk tree.  Every other component has its ordinary boundary output
at the node pointing toward that root.

Let $\eta$ be a closed smooth $(n-2)$-form on $D$.  Pulling back $\eta$
at the interior mark and the inputs at the boundary marks, and
integrating, gives functionals
\begin{equation}\label{curv:bulk-derivative}
 p_k^\eta(x_1,\ldots,x_k)
 =\sum_\beta T^{\omega(\beta)}p_{k,\beta}^\eta(x_1,\ldots,x_k),
 \qquad k\geq1.
\end{equation}
We include all words of diagonal and transverse branch sectors, with
their capping orientation lines.  The signs use the suspended degrees
$|x|'=|x|-1$.  Thus a cyclic permutation of the inputs has its usual
Koszul sign in these suspended degrees.  On a continuous-lift sector
the virtual dimension of the underlying disk problem is
\begin{equation}\label{curv:unmarked-dimension}
 \operatorname{vdim}\mathcal P_{k,\beta}=n+k-1.
\end{equation}
The capped index gives the corresponding formula on a word containing
branch sectors.  The operation is zero unless its integrand has top
degree.  Spaces with $k=0$ are used for Stokes; the unstable constant
disk with just the interior mark is excluded.

\begin{lemma}[Interior-rooted virtual integration]\label{curv:stokes}
Relative to the fixed cyclic boundary-rooted operations, these spaces
admit oriented virtual integrations through every finite cutoff.
Their normalized boundary products have an interior-rooted parent
and an ordinary boundary-rooted child, matched either over the source
diagonal or at reciprocal branch sectors.  The integrations are
cyclic in their external inputs and obey Stokes and composition.
In particular,
\begin{equation}\label{curv:pairing-identity}
 p_1^\eta(m_0)=0,
 \qquad
 p_{1,0}^\eta(c)=\pm\int_L c\wedge i^*\eta
 \quad(c\in\Omega^2(L)).
\end{equation}
The sign in the second identity is fixed and nonzero.  The ordinary
disk operations retain their chosen values through the cutoff.
\end{lemma}

\begin{proof}
We construct the interior-rooted spaces with all immersed boundary
labels.  Their local boundary problems are the ordinary diagonal and
transverse-jump problems of \cite[Sections~4--5]{AkahoJoyce}; the
relative construction below specifies their organization when no
boundary output is chosen.  It uses the mixed parent--child
organization of \cite[Propositions~16.3 and~16.14]{AFOOO2026}, retaining
reciprocal internal switches even when the external word is empty.
The parent is always additional interior-rooted data.  Its
identification with a cyclic boundary-rooted operation is neither
required nor imposed.

\emph{Compactification and normalization.}
Keep the interior mark in immersed stable-map compactness.  At a
boundary node the two lifts either agree in $L$ or end at distinct
preimages $p,q$ of a double value.  In the latter case the two
normalized flags have the reciprocal labels $(p,q)$ and $(q,p)$.
These flags disappear from the external word on gluing; they remain
part of the normalized stratum.  In particular, an empty external
word does not restrict its compactification to the diagonal sector.

There are no nonconstant spheres.  A connected tree of constant
spheres attached at one interior node and carrying at most the one
interior mark cannot be stable: it has fewer than three external
special points in total.  The degeneration in which the entire
boundary collapses would retain a nonconstant sphere, and is also
absent.  All remaining strata are disk trees.  Constant stable
vertices are retained, including the disk with the interior mark
and one boundary node.  The constant disk with no boundary input or
node and only the interior mark is unstable and gives no stratum.

Let $v_*$ be the vertex carrying the interior mark.  Orient every
edge toward $v_*$.  A vertex other than $v_*$ then has exactly one
outgoing flag, which is its ordinary disk output; every remaining
flag is an input.  A normalized corner is the product of the moduli
problems on these vertices, with the diagonal matching equations
and reciprocal-label matches along its edges.  For any collection
of edges, cutting them in any order produces this same oriented
tree and the same factors.  The root of each separated component
is either $v_*$ or its unique edge toward $v_*$.  This rule also
describes intermediate faces, so the identifications on iterated
normalized corners agree.

Stabilization respects these identifications.  Forgetting auxiliary
marks contracts only constant unstable components.  Contracting a
non-root such component carries the unique path toward $v_*$ to the
unique path toward $v_*$; hence the output flag on every surviving
ordinary component is unchanged.  A constant component carrying the
interior mark and an attaching boundary node is stable and cannot
be contracted.  If it has no such node, it is the excluded constant
zero-input root.  An unstable constant bivalent disk has either two
diagonal nodes on the same sheet or two reciprocal switch nodes:
its two constant boundary arcs rule out exactly one switch.  Its
contraction preserves the labels of the surviving flags.  The
two-switch constant strip is the identity under the dual cap-line
identification fixed by determinant gluing; it is not a new
orientation choice \cite[Sections~4.3 and~5]{AkahoJoyce}.
Thus a contraction performs the same pairing
of the ordered cap lines whichever order of normalization is used.
For a chain of these components, the neck lengths concatenate as
$T=T_1+\cdots+T_j$.  This is associative.  Forgetting is interpreted
in the admissible category of \cite[Section~14.3]{AFOOO2026}:
its pullbacks of admissible obstruction and CF data are admissible
and smooth in the gluing profiles used for integration.  The
forgetting map itself need not be smooth in ordinary corner
coordinates.  This admissible compatibility gives compatibility
with stabilization as well as with cutting a fixed tree.

These constructions are equivariant under isomorphisms of decorated
trees.  Such an isomorphism fixes the interior mark and preserves
the directed paths to it, so it never changes an ordinary child's
output.  Relabeling external inputs permutes their positions without
rerooting any ordinary child.  The same statement holds for an
automorphism permuting several isomorphic subtrees.  Normalized
corners retain their ordering of the cut edges; changing that
ordering permutes the outward normal lines with their determinant
sign.  Their subsequent finite quotients therefore have the required
corner identifications and automorphism weights.

\emph{Relative charts and smoothing.}
At a smooth nonconstant root choose auxiliary marked points on
compact cores, in free orbits of its finite automorphism group,
and impose transverse incidence conditions.  Choose obstruction
sections supported away from the distinguished mark, the nodes,
and the strip ends.  Their span can make the augmented
Cauchy--Riemann operator surjective: a cokernel vector annihilating
all such sections would vanish on an open core and then vanish by
unique continuation.  Taking the span of the stabilizer translates
makes this an invariant finite-dimensional obstruction space.
If effective charts are used, include vectors $D_u\xi$ for map
variations $\xi$ supported on disjoint free stabilizer-orbit balls
in the core.  These variations lie in the augmented kernel and
carry a faithful stabilizer representation modulo infinitesimal
reparametrizations; they make the thickened chart effective.
No submersivity is requested at the interior evaluation, which is
used only to pull back $\eta$.

At a broken root use the direct sum of the obstruction data on its
normalized components.  Ordinary vertices carry the fixed ordinary
data.  At a diagonal edge the child output evaluation is submersive,
so component surjectivity remains surjectivity after imposing the
matching equation in $L$.  All matching equations can be solved
recursively from the root toward the leaves: having chosen the
parent variation, choose the child's variation with the required
output value and continue with its children.  Joint submersivity
of all boundary evaluations on one component is not needed.
At a reciprocal switch the matching
space is a point; the operator has precisely the transverse
fixed-jump asymptotic condition of the immersed disk charts.
The linear and nonlinear gluing problems are therefore those of
the ordinary immersed construction, with the interior-mark
coordinate added on the root component.  A constant root
stabilized by that mark is kept as a marked component in this
description.

Transport obstruction sections from the compact cores into the
plumbed domains.  Exponential gluing estimates give charts through
the smoothing parameters, with the corrected smooth gluing
profiles cited above.  Coordinate changes are obtained from
componentwise obstruction data.  We specify their relative choice,
so that compatibility is more than a choice of unrelated local
thickenings.  Use the quasi-component construction of
\cite[Conditions~8.12--8.15, Lemmas~8.19--8.20, and
Definition~8.21]{FOOOChartsII}.  A quasi-component records a
transported compact-core obstruction choice and its auxiliary
stabilization data.  On a boundary tree take the componentwise
union of the already chosen indices: those of its interior root
and the fixed indices of its ordinary vertices.  The directed-tree
identifications show that this prescription agrees on every
iterated boundary and stabilizer quotient.

Extend the corresponding open family of indices inside a proper
family over a sufficiently small boundary neighborhood, as in
\cite[Proposition~8.18]{FOOOChartsII}.  Complete its finite cover
by new compact-core choices on the root away from a smaller
neighborhood, and saturate the choices under the finite
automorphism groups.  Choose the new compact-core spaces generically
transverse to the finitely many inherited spaces and their
stabilizer translates, as in the quasi-component construction.
Surjectivity and these direct-sum conditions are then open;
shrinking the neighborhoods preserves them.
Properness gives the semicontinuity inclusion of obstruction
spaces under changes of the central map, as in
\cite[Proposition~8.17]{FOOOChartsII}.  These inclusions and
forgetting the auxiliary marks define the coordinate changes;
they satisfy the cocycle conditions because the index inclusions
compose and the admissible forgetting operations commute with
normalization.  The original boundary families are unchanged on
the smaller collar.  This is the componentwise proof of
\cite[Theorems~5.3 and~5.4]{FOOOChartsII}, with its additional root
data specified relatively.  The local diagonal and fixed-jump
operators and the augmented matching surjectivity verified above
are exactly what that construction uses.  The fixed ordinary
factors are retained.

\emph{Finite extension of perturbations.}
Let $\delta>0$ be a lower bound for positive component energies,
and choose $0<\varepsilon<\delta$.  Assign cost
$E+\varepsilon(k-1)$ to an ordinary vertex of energy $E$ with $k$
inputs and cost $E+\varepsilon k$ to an interior-rooted vertex.
These costs add under boundary splitting.  Every ordinary stable
vertex has positive cost: the zero-energy unary strip is unstable
and represents the separately fixed de Rham differential, and a
zero-energy stable ordinary vertex has at least two inputs.  The
constant one-input interior root has cost $\varepsilon$; the
constant zero-input root is excluded.  Each nontrivial breaking
therefore has lower-cost factors.  The locally finite energy set,
the arity cutoff, and stability give a finite collection of
problems and normalized corner types, after adjoining all required
faces.  This includes constant switching polygons; their number of
inputs is controlled by arity, not by their zero energy.

Induct on this finite order.  On every normalized boundary prescribe
the product of the already constructed interior-rooted parent
perturbation and the fixed ordinary child perturbation.  The
prescriptions agree on higher corners by the directed-tree and
stabilization identifications.  Diagonal fiber products are defined
because the child output is strongly submersive for its fixed
CF-perturbation; reciprocal switches are products over points.
First extend these compatible normalized-boundary products to a
corner neighborhood by outer collaring, as in the collar step of
\cite[Proposition~14.8 and Remark~14.7]{AFOOO2026}.  The required
agreement on iterated normalized boundaries is precisely the
agreement verified above, including their stabilizer quotients.
Apply \cite[Proposition~14.5]{AFOOO2026} with the map to a point
and with the closed smaller corner neighborhood as its prescribed
set.  It supplies a compatible thickening and extends the
CF-perturbation over the new root space while retaining it near
that set.  Composition on the prescribed collar still uses the
original ordinary child perturbations.  There is no requirement
to forget a boundary input on the new root space and no new
submersion condition at its interior mark.

For cyclic relabeling and finite stabilizers, take invariant
obstruction spaces and include the finite-group representations in
the auxiliary parameter spaces of the CF-perturbations.  The
prescribed corner products are equivariant, so the same relative
extension is equivariant.  This uses equivariant families, not
single-valued equivariant transverse sections.  Finite quotients
are integrated with their ordinary stabilizer weights.  Stokes and
composition now apply to every prescribed product.  The identical
construction over a parameter interval, relative to its endpoints,
supplies finite-cut comparisons whenever the ordinary choices vary.

\emph{Orientations and the zero-input identity.}
Use the spin determinant orientations and the graded cap lines at
all transverse flags.  At a reciprocal node the two cap lines are
paired in their ordered duality; at a diagonal node they are
replaced by the matching determinant of $TL$.  Determinant gluing
is associative.  Permuting normalized smoothing normals or cap
lines gives exactly their Koszul signs, including at an
automorphism quotient.  Stabilizing marks together with their
incidence constraints contribute canceling orientation factors.
These are the orientation conventions for ordinary immersed disk
composition; adding the even-dimensional interior-mark coordinate
does not alter its cap-line rule.

For $k=0$, pull back the closed $(n-2)$-form $\eta$ to the virtual
$(n-1)$-space.  Each boundary product has a single ordinary
curvature child and an interior-rooted parent with one input.
Composition first identifies its contribution with
$p_1^\eta(m_0)$, with the common Stokes sign.  This includes every
reciprocal-switch product.  Only after composition do we use the
grading: a curvature output has degree two, while each branch
generator has odd degree, so those child outputs vanish.
There is no sphere or boundary-collapse term.  Stokes gives the
first identity in \eqref{curv:pairing-identity}.

At energy zero the one-input diagonal root is the regular constant
disk space $L$.  Its one boundary mark and one interior mark fix
the domain; its evaluations are the identity and $i$.  Its integral
is the second identity in \eqref{curv:pairing-identity}.  The
boundary-rooted operations were never changed in this construction.
\end{proof}

\subsection{Leading obstructions and canonical models}

For a degree-one cochain $b$ of positive valuation, put
\[
 m_0^b=\sum_{k\geq0}m_k(b,\ldots,b),
 \qquad
 R_\eta^b(c)=\sum_{r\geq0}p_{r+1}^\eta(c,b,\ldots,b).
\]
These sums converge in the energy filtration.  The cochain $b$ may
have both diagonal and degree-one branch components.  Through a
fixed cutoff only boundedly many insertions occur.

\begin{lemma}[Boundary deformations]\label{curv:deformation}
For the fixed ordinary operations and the relative integrations of
Lemma~\ref{curv:stokes},
\begin{equation}\label{curv:deformed-pairing}
 R_\eta^b(m_0^b)=0.
\end{equation}
If $m_0^b$ has a first nonzero energy coefficient, that coefficient
is a closed exact two-form on $L$.
\end{lemma}

\begin{proof}
We first establish the identity using the complete virtual boundary
products.  Apply Stokes to the interior-rooted spaces with $k$
external inputs all equal to $b$.  A boundary child carrying $s$
consecutive inputs has output $m_s(b,\ldots,b)$, and the root carries
the remaining $r=k-s$ inputs and this output.  Differentiating an
input form supplies the zero-energy differential term in $m_1$.
There is no differential of $\eta$ and no boundary output
differential.

The orientation signs are the suspended insertion signs.  At a
diagonal matching node this is the ordinary disk Stokes calculation,
with the even interior mark in place of the root's boundary-output
coordinate.  External branch cap lines are moved with their total
shifted degrees in the same calculation.  Since every input $b$
has suspended degree zero, all insertion signs are $+1$; rotating
the degree-two curvature output past these inputs also has sign
$+1$.  This applies to degree-one point sectors as well as to
one-forms.  Reciprocal child outputs are retained in the product
formula, and then vanish because they would have degree two in an
odd-degree point sector.

For $k>0$ there are $k$ cyclic positions for each contribution.
For $0<s<k$ they are the starts of the child block; for $s=0$ they
are the $k$ gaps; and for $s=k$ they are the $k$ gaps at the child's
attaching node, or equivalently its rooted linear orders.  The
input differentials also have $k$ positions.  These are counts on
normalized labeled strata, with stabilizer weights retained; we
do not quotient the spaces by the external cyclic action.
Dividing Stokes by $k$ thus gives
\[
 \sum_{r+s=k}
 p_{r+1}^\eta\bigl(m_s(b,\ldots,b),b,\ldots,b\bigr)=0.
\]
The separate $k=0$ identity is $p_1^\eta(m_0)=0$, proved above.
Summing over $k$ proves \eqref{curv:deformed-pairing}.  Through any
cutoff the sum is finite, so all these operations and divisions are
legitimate.

Let $c_E$ be the first nonzero coefficient of $m_0^b$.  Grading
places it in $\Omega^2(L)$, since no branch sector has degree two.
The leading curved $A_\infty$ equation is $dc_E=0$.  Every term
of $R_\eta^b$ involving a $b$ insertion or positive disk energy has
positive valuation.  Its remaining term is
$c\mapsto\pm\int_L c\wedge i^*\eta$.  The coefficient of energy
$E$ in \eqref{curv:deformed-pairing} is therefore
\begin{equation}\label{curv:ambient-annihilation}
 \int_L c_E\wedge i^*\eta=0.
\end{equation}
Choose closed representatives of a basis of $H^{n-2}(D;\C)$.
Equation~\eqref{curv:ambient-annihilation} says that
$i_*([c_E]\cap[L])=0$.  The injection in
\eqref{curv:topological-input} remains injective over $\C$.
Poincar\'e duality on $L$ now gives $[c_E]=0$.
\end{proof}

\begin{proof}[Proof of Proposition~\ref{prop:curvature}]
First consider the cyclic canonical model on
\eqref{curv:canonical-space}.  Suppose its curvature has a first
nonzero coefficient $c_E^{\mathrm{can}}$.  The local finiteness of
the energy filtration ensures that a first coefficient exists if the
curvature is nonzero.  Filtered transfer from this model to the
geometric forms model includes an arity-zero component $F_0$ of
positive valuation.  Set $b=F_0$.  The zero-input functor equation is
\begin{equation}\label{curv:transfer-equation}
 m_0^{\mathrm{forms},b}
       =F_1(m_0^{\mathrm{can}}).
\end{equation}
Its leading term is
$F_{1,0}(c_E^{\mathrm{can}})$, where $F_{1,0}$ induces the ordinary
cohomology identification.  Since the canonical degree-two summand
is $H^2(L;\C)$, this is a nonzero cohomology class.  But
Lemma~\ref{curv:deformation} says that the first geometric curvature
coefficient is exact, a contradiction.

The same argument applies to a canonical structure obtained by
promotion from finite cutoffs.  Choose a geometric cutoff strictly above $E$, and compose transfer
with the cutoff pseudo-isotopy.  Apply Lemma~\ref{curv:stokes}
relative to the ordinary operations on this geometric representative.  Its arity-zero
translation still has positive valuation, and its leading unary map
still induces an isomorphism on ordinary cohomology.  Thus
\eqref{curv:transfer-equation} through that cutoff gives the same
contradiction.  All these comparisons use the energy/class filtration
of the fixed $J_0$ disk system.  Lemma~\ref{curv:energy} supplies the common gap for the ordinary
operations and their comparisons.  The interior functionals only
need to be constructed on the chosen geometric representative; no
identification of them across the pseudo-isotopy is used.  Terms containing positive-valuation cochains require
only bounded arity below a fixed energy.  Equivalently, one may carry
out the argument in filtered $A_{n',k'}$ models, with both cutoffs
larger than those needed to detect $c_E^{\mathrm{can}}$.
The canonical curvature is therefore zero.

For an embedded test copy there are no positive-energy self-disks
at $J_0$, so its geometric curvature is zero.  Keep $J_0$ fixed in
the pseudo-isotopies and choose transfers with zero arity-zero
translation on these objects.  The self-disk contributions, including
the arity-zero homotopy terms, are empty, so promotion preserves this
choice.  On a strictly increasing word all morphisms are isolated
intersection generators.  Energy-zero transfer is the identity on
these point complexes, and its restriction needs no operations with
repeated labels.  The operations on this sector therefore remain the
geometric disk operations up to the stated cutoff pseudo-isotopies.

Homotopy units give cohomological units on the resulting uncurved
objects.  For $L$ this can also be seen by transferring the zero
canonical bounding cochain to the forms model: its image is the
positive cochain $F_0$, which is a bounding cochain by
\eqref{curv:transfer-equation}.  Deform the homotopy units there and
use the deformed filtered quasi-isomorphism.  No strict chain unit
is required for this step.  All operations are graded, so the
cohomological unit has degree zero.

After extending to $\Lambda$, the canonical morphism complexes are
finite-dimensional.  The self complex of $L$ has degree-two chain
space exactly $H^2(L;\C)\otimes\Lambda$, because all branch
generators have odd degree.  Its degree-two cohomology is a
subquotient of that space, proving \eqref{curv:degree-two-bound}.
Finally, the nondegenerate cyclic pairing gives duality with shift
$n$ on morphism cohomology.  Its extension to twisted complexes and
duality on retracts gives the same pairing on the split-closed
triangulated hull, whose Serre functor is consequently $[n]$.
\end{proof}
